\section{Latent Generalization Gap：为什么反向关系最先暴露问题}

最清楚的入口是 reversal curse。若模型学到“Daphne Barrington 导演了 A Journey Through Time”，人类通常会认为“谁导演了这部电影”只是同一关系的反向查询；但 fine-tuned LLM 往往只能沿训练语句的方向回答。奇怪的是，把同一句新事实写进聊天上下文，模型却能轻松反向使用。

\subsection{Reversal curse 与聊天中的反例}

设二元关系 $R(a,b)$ 表示“$a$ 与 $b$ 具有训练中给定的方向关系”。反向问题要求模型调用
$$
R(a,b)\Longrightarrow R^{-1}(b,a).
$$
这里：
\begin{itemize}
  \item $R$ 是原始关系，如“导演了”或“包含”；
  \item $R^{-1}$ 是查询方向相反的关系；
  \item 逻辑信息并未增加，变化的是访问方向与输出形式。
\end{itemize}

\lecturefigure{slide-007.jpg}{Fine-tuning 后出现的 reversal curse}{Stanford CS25 V6 Lecture 06 官方 slide p. 7}

\teachervoice{00:03:30--00:04:21，老师说真正令人惊讶的不是 reversal curse 本身，而是“在 chat 里这明明很容易”。他之前甚至把 reversal 当作弱模型也能完成的 easy baseline，这个反差促成了参数与上下文的直接比较。}

\lecturefigure{slide-008.jpg}{上下文中的模型能够直接完成同一反向关系}{Stanford CS25 V6 Lecture 06 官方 slide p. 8}

\begin{importantbox}{核心矛盾不是缺少事实，而是事实的调用方向}
Fine-tuned 模型已经在正向问题上证明它记住了事实；失败发生在目标函数没有直接训练过的反向访问上。ICL 的成功则说明同一架构在信息仍位于 context 时，具备完成这种变换的计算能力。
\end{importantbox}

\subsection{“ICL 类似梯度下降”为什么不能结束讨论}

一些理论工作在特定线性回归、优化最优或构造性设定中发现，Transformer 的前向计算可以实现类似 gradient descent 的更新。因此 slide 10 故意提出一个反问：如果 ICL 在机制上近似某种优化，为什么真正更新参数后反而更不灵活？答案是，\textbf{局部算法相似}不等于\textbf{最终表示、压缩与访问接口相同}。

\lecturefigure{slide-009.jpg}{研究问题：参数与上下文是否诱导不同的泛化}{Stanford CS25 V6 Lecture 06 官方 slide p. 9}

\lecturefigure{slide-010.jpg}{ICL 与 gradient descent 的受限类比}{Stanford CS25 V6 Lecture 06 官方 slide p. 10}

\begin{warningbox}{ICL $\simeq$ gradient descent 不是全局等价定理}
这些结果通常依赖任务分布、模型构造和最优算法假设。即使前向激活执行了一步类似优化的计算，也不能推出：它与长期训练后得到相同参数表示、相同信息瓶颈、相同可逆性，或相同的 out-of-distribution generalization。
\end{warningbox}

\subsection{同一数据：整包放进 context，或完整 fine-tune}

为了避免“ICL 只看到了被人工挑选的相关事实”这一混淆，实验把整个新数据集放进长上下文，而不是只给 few-shot 示例。另一组模型则在同一数据上 fine-tune。测试同时包含显式方向和训练中未直接出现的 latent query，从而区分记忆、检索与关系变换。

\lecturefigure{slide-011.jpg}{Fine-tuning 与整数据集 ICL 的对照协议}{Stanford CS25 V6 Lecture 06 官方 slide p. 11}

可定义一个简化的 latent generalization gap：
$$
\Delta_{\text{latent}}=
\operatorname{Acc}_{\text{ICL}}(T_{\text{latent}})
-\operatorname{Acc}_{\text{FT}}(T_{\text{latent}}).
$$
其中：
\begin{itemize}
  \item $T_{\text{latent}}$ 是答案没有以目标形式直接出现在训练数据里的测试；
  \item $\operatorname{Acc}_{\text{ICL}}$ 是整数据集位于上下文时的准确率；
  \item $\operatorname{Acc}_{\text{FT}}$ 是同数据用于 fine-tuning 后的准确率；
  \item $\Delta_{\text{latent}}>0$ 只描述当前实验差异，不是所有任务上的普遍常数。
\end{itemize}

\lecturefigure{slide-012.jpg}{Reversal 数据集上 ICL 与 fine-tuning 的巨大差距}{Stanford CS25 V6 Lecture 06 官方 slide p. 12}

读这张柱状图时先分清 forward 与 reversal。Fine-tuned 模型在 forward 方向并非完全失忆，它的问题集中在 relation reversal；ICL 即使承载整套数据仍能显著恢复反向查询。图支持“访问与重组方式不同”，但不能单独说明参数里究竟缺少了哪一个可解释变量。

\begin{knowledgebox}{读图：先比较同一测试方向，再比较学习路径}
不要把蓝色、灰色与黄色柱简单理解成三个模型强弱。第一步先看 forward 条件，确认各方法确实吸收了显式事实；第二步再看 reversal 条件的掉点；第三步才比较 ICL 与 fine-tuning 的相对差异。只有这样才能把“没有学到事实”与“没有学会反向调用”分开。
\end{knowledgebox}

\lecturefigure{slide-013.jpg}{Syllogism 测试同样显示 context 优势}{Stanford CS25 V6 Lecture 06 官方 slide p. 13}

Syllogism 要求从“All zamp are snaff”与“No snaff are plusk”推出“No zamp are plusk”。这不是简单倒置一个短语，而是组合两个显式命题得到一个未写出的结论。ICL 的优势因此从 directional reversal 扩展到 relation composition。

\lecturefigure{slide-014.jpg}{第一阶段结论：ICL 对若干 latent test 的泛化更强}{Stanford CS25 V6 Lecture 06 官方 slide p. 14}

\teachervoice{00:08:55--00:12:15，老师特别强调：ICL 不是只拿到一条被检索好的事实，而是拿到整个数据集；不同模型和任务上的差距大小会变化，但 reversal 与 syllogism 的定性模式反复出现。}

\subsection{为什么互联网文本会训练 ICL，却没有训练出同样的参数行为}

互联网文本天然包含“先给背景、后问问题”“先定义实体、后进行指代与组合”的局部结构，因此语言模型在预训练中频繁练习从 context 临时建立关系。Slide 15 的难题在于：如果这种能力已经存在，为什么用新的关系 fine-tune 后，模型不把同样的变换规则迁移到参数知识上？这说明学习规则不仅决定存了什么，也决定知识以何种接口可用。

\lecturefigure{slide-015.jpg}{上下文结构常见，但 fine-tuning 后仍未学会相同泛化}{Stanford CS25 V6 Lecture 06 官方 slide p. 15}

\begin{knowledgebox}{一种工作解释：skill 与 content 分离}
预训练可能在参数中形成了“如何读取 context 并组合其中关系”的 skill；新事实留在 context 时，这个 skill 可直接作用于它。Fine-tuning 把 content 写入参数后，模型需要另一个接口把参数知识重新暴露给该 skill。若表示没有形成稳定的可调用方向，已有推理能力也无从使用。
\end{knowledgebox}

\subsection{从头训练仍然失败：不是旧模型的适配偶然}

一个自然反驳是：预训练模型早已形成固定表示，少量 fine-tuning 才导致 reversal failure；如果从头训练，也许网络会学会对称关系。课程用约 20M 参数的小模型、20,000 个合成关系和 held-out reversals 检验这一点。合成 token 没有英语语义，模型也不能借助世界知识猜答案。

\lecturefigure{slide-016.jpg}{假设：从头预训练也许能学会 latent relation}{Stanford CS25 V6 Lecture 06 官方 slide p. 16}

\lecturefigure{slide-017.jpg}{From-scratch reversal 实验：训练方向与 held-out 反向关系}{Stanford CS25 V6 Lecture 06 官方 slide p. 17}

若训练集中对大多数关系同时提供正反方向，只对一小部分关系隐藏 reverse，可把测试写成
$$
D_{\text{train}}=D_{\text{forward}}\cup D_{\text{reverse,seen}},
\qquad
T=D_{\text{reverse,heldout}}.
$$
其中模型已经看到“关系可以双向表达”的总体规律，却仍需把这条规律迁移到具体 held-out 实体。结果表明，parametric learning 对具体实例的 latent reversal 仍然薄弱，而把对应定义放入 context 可以完成编码。

\teachervoice{00:13:24--00:14:02，老师回答课堂提问时说明：from-scratch 条件里模型不知道英语，输入只是合成 token。这一控制排除了“模型根据名字、常识或预训练语义补全答案”的解释。}

\subsection{Codebook：从反向查询扩展到翻译式组合}

Codebook 任务先定义索引到字符的映射，再给若干编码示例。显式测试只要求使用训练中直接练过的索引；latent test 则要求把定义和编码规则组合起来，处理训练编码样本中没出现过的索引。它更接近临时语言、符号翻译和工具协议，而不只是自然语言句子的倒序。与 reversal 相比，这个任务把“关系方向”升级成“先读取定义、再执行程序”，因而能检查模型是否把定义当作可复用规则，而非只拟合见过的输入输出对。

\lecturefigure{slide-018.jpg}{Codebook 的 trained indices 与 latent indices}{Stanford CS25 V6 Lecture 06 官方 slide p. 18}

若 codebook 为 $c:k\mapsto v$，字符串编码可写为
$$
\operatorname{Encode}_{c}(s_1\cdots s_n)
=\bigl(c^{-1}(s_1),\ldots,c^{-1}(s_n)\bigr).
$$
这里：
\begin{itemize}
  \item $c$ 是索引到符号的映射；
  \item $c^{-1}$ 是从符号找回索引的逆映射；
  \item latent test 要求把“定义映射”与“执行编码”组合，而不是复述训练样例。
\end{itemize}

\lecturefigure{slide-019.jpg}{即使从头训练，若干 latent structure 仍不稳定泛化}{Stanford CS25 V6 Lecture 06 官方 slide p. 19}

\begin{importantbox}{从头训练实验排除了什么，没排除什么}
它排除了“只是旧模型不适应新事实”的简单解释，也说明训练分布中存在一般结构并不保证每个实例都以可重组方式编码。它没有证明所有架构、优化器、数据规模或学习规则必然失败；相反，它把“怎样学习可调用的 latent structure”变成后续算法问题。
\end{importantbox}

\subsection{显式内容、latent information 与可用信息}

课程随后给出最重要的概念压缩。一次经验的 explicit content 是表面直接陈述；latent information 是由这些陈述可推出、但没有以未来问题形式写出的关系；usable information 则还要再加一层条件：模型能否在当前目标下把 latent relation 变成正确输出。

\lecturefigure{slide-020.jpg}{Reversal、multi-hop、alternative goals 与 cross-lingual 的 latent structure}{Stanford CS25 V6 Lecture 06 官方 slide p. 20}

\begin{knowledgebox}{读图：四个例子共享同一个信息结构}
读图时不要逐个记任务名，而要寻找共同模式：左侧或上方给出 explicit experience，右侧或下方测试一个未被直接写出的目标。Reversal 改变方向，multi-hop 增加组合深度，alternative goal 改变行动目标，cross-lingual 改变表达语言；它们都在测试经验能否被重新解释。
\end{knowledgebox}

\begin{knowledgebox}{First-use glossary：latent learning}
\begin{tabularx}{\textwidth}{L{0.23\textwidth}X X}
\toprule
术语 & 含义 & 例子 \\
\midrule
Explicit content & 经验中直接出现的命题或目标 & “X is Y's parent” \\
Latent information & 由显式内容蕴含、但未按测试形式写出的信息 & “Y 的父母是谁？”可由关系反向得到 \\
Latent generalization & 在新方向、新组合或新目标下使用 latent information & 从两条 parent 关系推出 grandchild \\
Usable information & 模型在当前计算路径中能够稳定访问并用于输出的信息 & 事实存在参数中，但必须能被目标条件激活 \\
\bottomrule
\end{tabularx}
\end{knowledgebox}

\lecturefigure{slide-021.jpg}{参数整合更广，上下文经历细节更丰富}{Stanford CS25 V6 Lecture 06 官方 slide p. 21}

这张图不能被读成“参数差、上下文好”。参数学习的优势是跨文档整合、压缩和摊销；context 的优势是把少量经历以相对完整的形式暴露给模型。真正的张力是 breadth versus flexibility：覆盖更多经验时，是否牺牲了按未来未知目标重新解释单次经历的能力。

\subsection{为什么现实数据会掩盖失败：co-occurrence shortcut}

在自然语料中，“eagle”“bird”“wing”经常共同出现。即使模型没有执行严格的 syllogism，也可能根据词共现答对“eagles have wings”。因此实践中的非零准确率不能自动证明 systematic reasoning；它可能来自有效但不同的统计路径。

\lecturefigure{slide-022.jpg}{词共现可绕过显式 reversal 或 syllogistic inference}{Stanford CS25 V6 Lecture 06 官方 slide p. 22}

可以把预测信号粗略分成
$$
\log p(y\mid x)=S_{\text{stat}}(x,y)+S_{\text{rel}}(x,y)+\varepsilon,
$$
其中：
\begin{itemize}
  \item $S_{\text{stat}}$ 表示词频、实体相似性与共现等统计信号；
  \item $S_{\text{rel}}$ 表示方向变换、组合与规则应用等关系信号；
  \item $\varepsilon$ 汇集未建模因素。该式是教学分解，不是论文声称的可辨识生成模型。
\end{itemize}

\begin{warningbox}{统计捷径不是“作弊”，但会改变证据含义}
统计规律通常提高真实世界的期望性能，不能因为它不是形式逻辑就删除。问题在于：若评测目标是验证 reversal 或 composition，数据中的共现线索会让模型绕过待测机制。受控实验需要消除捷径；部署系统则应同时利用统计与结构推理。
\end{warningbox}

\teachervoice{00:18:39--00:21:26，老师先说 co-occurrence 会掩盖 latent failure，随后又明确纠正一种过度解读：统计结构本身非常有用，课程并不是主张模型只能靠纯逻辑。}

\lecturefigure{slide-023.jpg}{没有关联线索时，怎样让模型学会 latent relation？}{Stanford CS25 V6 Lecture 06 官方 slide p. 23}

\subsection{本章小结}

Reversal、syllogism 与 codebook 共同表明：模型能记住 explicit content，不代表它能把同一经验转换成任意未来目标需要的形式。ICL 在这些受控任务上更灵活，但优势并非普遍定理；现实语料中的共现、语义与检索技能会改变结果。下一步不是简单选择 context，而是把 context 的灵活计算迁移到可扩展系统中。

